\subsection{Materia necesaria}

\begin{definicion}{Sistema \(G/G/c\)}
Cuando hay llegadas generales, servicios generales y \(c\) servidores, se puede aproximar la espera con Sakasegawa. Primero se calcula utilizacion:
\[
\rho=\frac{\lambda}{c\mu}.
\]
Si \(\rho\ge 1\), el sistema no es estable.
\end{definicion}

\begin{formulas}{Sakasegawa}
\[
W_q\approx
\left(\frac{CV_a^2+CV_s^2}{2}\right)
\frac{\rho^{\sqrt{2c+2}-1}}{c(1-\rho)\mu}.
\]
\end{formulas}

\begin{algoritmo}{Asignacion de capacidad}
\begin{enumerate}
    \item Convertir demanda de cada bloque a tasa horaria \(\lambda_b\).
    \item Calcular \(\mu\) por taxi.
    \item Calcular \(\rho_b\); si supera 1, falta capacidad.
    \item Para una meta de utilizacion, despejar \(c_b\ge \lambda_b/(\rho^*\mu)\).
\end{enumerate}
\end{algoritmo}
